\section{Conclusion}
\label{sec:conclusion}

Static checks on agent workflows are only as credible as the behavioral model
they run on, and recovering that model from public agent code---not the check
algorithms---is the binding constraint. On a cross-framework corpus of $922$
extracted file-level graphs from $252$ public GitHub repositories, extraction
error dominates check error and can reverse both apparent precision and
apparent prevalence; the DSL and monitor-pruning layer are supporting
infrastructure for that finding, not the headline.

The estimates are textured, not a null. Read as distinct estimands rather than
a pooled rate, LLM-assisted exploratory estimates put confirmed structural
defects at $0.84\%$ and human-gate policy violations at $2.5\%$ of the
validation sample (Table~\ref{tab:estimands}), with the composite
post-stratifying to $2.2\%$; three of four observed cases occur in
application-like files, too few to establish a stratum difference (Fisher exact
$p{=}0.10$). Yet flag-level precision stays low (workflow-level PPV $1.8\%$),
none of the structural flags on extracted graphs was a confirmed defect, and
the diagnosis holds across every cut---edge recall $0.68$ on real code, $0.36$
on implicit-topology ADK. The static machinery still keeps a defect-independent
role: under certified event-trace containment the graph\,$\times$\,DFA product
soundly prunes $237$ of $270$ curated monitor instances ($88\%$;
Table~\ref{tab:pruning}), a monitor-selection layer in a static-plus-runtime
safety stack rather than a scale result.

\paragraph{Limitations.}
The population is public GitHub, not production systems; two thirds of it is
tutorial-grade, which we census and stratify but cannot fully correct. The unit
of analysis is one extracted file-level graph, not an executable workflow root;
a root-level unit is future work. Reference graphs and triage labels are
LLM reconstructions with adversarial verification, not human annotations;
confidence-stratified sensitivity checks and released per-label justifications
mitigate but do not eliminate this, and a two-human independent validation is
staged (protocol and samples prepared) as the immediate next step. Prevalence
CIs are wide at $n{=}119$; the estimates are best read as ``rare, but real'',
not as precise rates. We read these results as a redirection---toward
higher-fidelity extraction and per-workflow policy instantiation.

\paragraph{Artifact.}
An anonymized artifact is submitted with this paper, not merely promised: the
tool, both corpora, the mining and validation pipeline, the LLM-reconstructed
reference graphs, triage labels with per-label justifications, and all
evaluation scripts, packaged for double-blind review under the MIT license.

\section*{Ethical Statement}
The study mines only \emph{public} GitHub repositories, records exact provenance
(repository at commit SHA plus file path), and redistributes only derived graph
abstractions, never third-party source; no personal data is processed. Confirmed
defects are reported without identifying maintainers, and we disclose them to
the affected repositories before publication. The intended impact is
defensive---the honest finding that such defects are rare in public code should
calibrate, not inflate, claims about static verification's role in agent
safety.
